\documentclass[11pt,a4paper,oneside]{book}

\usepackage[spanish,es-tabla]{babel}
\usepackage[utf8]{inputenc}
\usepackage[T1]{fontenc}
\usepackage{lmodern}
\usepackage{microtype}
\usepackage[a4paper,margin=2.7cm]{geometry}
\usepackage{xcolor}
\usepackage{amsmath,amssymb,mathtools}
\usepackage{booktabs}
\usepackage{array}
\usepackage{longtable}
\usepackage{enumitem}
\usepackage{hyperref}
\usepackage{listings}
\usepackage{fancyhdr}
\usepackage{tcolorbox}
\usepackage{multicol}

\definecolor{azuluc}{HTML}{1F4E79}
\definecolor{grisclaro}{HTML}{F4F6F8}
\definecolor{gris}{HTML}{586069}
\definecolor{verde}{HTML}{2E7D32}
\definecolor{rojo}{HTML}{A94442}

\hypersetup{
  colorlinks=true,
  linkcolor=azuluc,
  citecolor=azuluc,
  urlcolor=azuluc,
  pdftitle={Manual de Laboratorio de Matemáticas},
  pdfauthor={Asignatura Laboratorio de Matemáticas}
}

\lstdefinestyle{rstyle}{
  language=R,
  basicstyle=\ttfamily\small,
  keywordstyle=\color{azuluc}\bfseries,
  commentstyle=\color{gris},
  stringstyle=\color{verde!70!black},
  showstringspaces=false,
  columns=fullflexible,
  frame=single,
  rulecolor=\color{gris!30},
  backgroundcolor=\color{grisclaro},
  breaklines=true,
  tabsize=2
}

\lstdefinestyle{texstyle}{
  language=[LaTeX]TeX,
  basicstyle=\ttfamily\small,
  keywordstyle=\color{azuluc}\bfseries,
  commentstyle=\color{gris},
  stringstyle=\color{verde!70!black},
  showstringspaces=false,
  columns=fullflexible,
  frame=single,
  rulecolor=\color{gris!30},
  backgroundcolor=\color{grisclaro},
  breaklines=true,
  tabsize=2
}

\lstdefinestyle{plainstyle}{
  basicstyle=\ttfamily\small,
  showstringspaces=false,
  columns=fullflexible,
  frame=single,
  rulecolor=\color{gris!30},
  backgroundcolor=\color{grisclaro},
  breaklines=true,
  tabsize=2
}

\lstset{
  literate=
    {á}{{\'a}}1 {é}{{\'e}}1 {í}{{\'i}}1 {ó}{{\'o}}1 {ú}{{\'u}}1
    {Á}{{\'A}}1 {É}{{\'E}}1 {Í}{{\'I}}1 {Ó}{{\'O}}1 {Ú}{{\'U}}1
    {ñ}{{\~n}}1 {Ñ}{{\~N}}1 {ü}{{\"u}}1 {Ü}{{\"U}}1
    {¿}{{?`}}1 {¡}{{!`}}1
}

\tcbset{
  colback=grisclaro,
  colframe=azuluc,
  boxrule=0.6pt,
  arc=2pt,
  left=8pt,
  right=8pt,
  top=6pt,
  bottom=6pt
}

\newtcolorbox{objetivos}{title=Objetivos de aprendizaje}
\newtcolorbox{entrega}{title=Entrega de la práctica,colframe=verde!60!black}
\newtcolorbox{aviso}{title=Observación,colframe=rojo!70!black}

\pagestyle{fancy}
\setlength{\headheight}{14pt}
\fancyhf{}
\fancyhead[L]{Laboratorio de Matemáticas}
\fancyhead[R]{\thepage}
\renewcommand{\headrulewidth}{0.4pt}

\setlist[itemize]{topsep=4pt,itemsep=2pt}
\setlist[enumerate]{topsep=4pt,itemsep=3pt}

\title{\Huge Manual de Laboratorio de Matemáticas\\[0.5em]
\Large Grado en Ingeniería Matemática}
\author{Borrador docente}
\date{Curso 2026/27}

\begin{document}

\frontmatter
\maketitle

\chapter*{Propósito del manual}
Este manual acompaña la asignatura \emph{Laboratorio de Matemáticas}, de primer curso del Grado en Ingeniería Matemática. La asignatura tiene 6 ECTS, 15 horas de clase magistral y 45 horas de prácticas presenciales. Su finalidad es introducir el uso de herramientas computacionales para resolver, documentar y comunicar problemas matemáticos.

El enfoque principal del manual es el ecosistema R: programación básica, cálculo numérico elemental, tratamiento de datos, representación gráfica y modelización estadística inicial. La edición científica se trabaja con LaTeX, y los cuadernos reproducibles se presentan mediante Markdown, R Markdown o Quarto.

\section*{Cómo usar este documento}
Cada bloque combina una parte conceptual breve con ejemplos reproducibles y una práctica de laboratorio. Las prácticas están pensadas para sesiones de 3 horas. El manual puede adaptarse a grupos con distinto nivel inicial ampliando o reduciendo los apartados marcados como extensión.

\section*{Convenciones}
\begin{itemize}
  \item El símbolo \texttt{>} identifica comandos ejecutados en la consola de R.
  \item Los fragmentos de código se pueden copiar en un script \texttt{.R}, salvo que se indique otra cosa.
  \item Las entregas prácticas deben incluir código, resultados y una explicación breve de las decisiones tomadas.
  \item Las soluciones no se valoran solo por dar un número correcto: importan la organización del proyecto, la claridad del código y la interpretación matemática.
\end{itemize}

\tableofcontents

\mainmatter

\chapter{Organización de la asignatura}

\section{Resultados de aprendizaje}
Al completar la asignatura, el estudiante debe ser capaz de:
\begin{itemize}
  \item utilizar R como entorno de trabajo para resolver problemas matemáticos sencillos;
  \item escribir programas con variables, funciones, condicionales, bucles y estructuras de datos;
  \item aplicar herramientas computacionales a cálculo, simulación, datos y estadística inicial;
  \item generar gráficos y comunicar resultados cuantitativos de forma clara;
  \item elaborar documentos científicos básicos en LaTeX;
  \item crear cuadernos reproducibles con Markdown, R Markdown o Quarto;
  \item organizar proyectos, scripts, datos y resultados de forma trazable.
\end{itemize}

\section{Distribución temporal}
La siguiente planificación distribuye las 15 horas magistrales y las 45 horas prácticas en 15 semanas. Cada semana contiene 1 hora magistral y 3 horas de laboratorio.

\begin{longtable}{p{0.08\textwidth}p{0.40\textwidth}p{0.40\textwidth}}
\toprule
Semana & Clase magistral & Práctica de laboratorio \\
\midrule
1 & Entorno de trabajo y proyectos & Instalación, RStudio/Posit y primer proyecto \\
2 & Objetos, tipos y expresiones & Cálculo elemental y scripts reproducibles \\
3 & Funciones y ayuda integrada & Funciones propias y documentación mínima \\
4 & Condicionales e iteraciones & Algoritmos elementales y comprobaciones \\
5 & Vectores y matrices & Operaciones vectorizadas y álgebra computacional \\
6 & Listas y tablas de datos & Importación, limpieza y exploración de datos \\
7 & Gráficos base y \texttt{ggplot2} & Visualización de funciones y datos \\
8 & Cálculo numérico elemental & Raíces, derivadas e integrales aproximadas \\
9 & Simulación & Experimentos aleatorios y estimación por Monte Carlo \\
10 & Estadística descriptiva & Resúmenes, distribuciones y comparación de grupos \\
11 & Modelización estadística inicial & Regresión lineal e interpretación de modelos \\
12 & LaTeX científico & Fórmulas, tablas, figuras y bibliografía \\
13 & Markdown y cuadernos reproducibles & Informe reproducible con R Markdown o Quarto \\
14 & Proyecto integrador & Desarrollo guiado de un informe computacional \\
15 & Presentación y revisión & Entrega, defensa breve y mejora final \\
\bottomrule
\end{longtable}

\section{Criterios de evaluación recomendados}
La guía docente permite combinar evaluación continua y examen final, cada uno con peso entre el 40\% y el 60\%. Una concreción coherente para esta asignatura es:

\begin{itemize}
  \item Evaluación continua, 50\%: prácticas semanales, informes reproducibles, ejercicios de programación y participación técnica en laboratorio.
  \item Examen final, 50\%: resolución individual de problemas de programación, cálculo, análisis de datos y edición científica.
\end{itemize}

En las prácticas se recomienda valorar:
\begin{itemize}
  \item corrección matemática y computacional;
  \item claridad y simplicidad del código;
  \item uso adecuado de funciones y estructuras de datos;
  \item reproducibilidad del trabajo;
  \item interpretación escrita de resultados;
  \item presentación formal del informe.
\end{itemize}

\chapter{Entorno de trabajo}

\section{R, RStudio y proyectos}
R es un lenguaje de programación orientado al cálculo estadístico, el análisis de datos y la representación gráfica. RStudio o Posit Desktop proporciona un entorno integrado para escribir scripts, ejecutar código, consultar ayuda, visualizar gráficos y organizar proyectos.

Un proyecto bien organizado evita pérdidas de información y facilita que otra persona pueda reproducir los resultados. Una estructura básica es:

\begin{lstlisting}[style=plainstyle]
laboratorio-matematicas/
  datos/
  scripts/
  informes/
  figuras/
  README.md
\end{lstlisting}

\section{Primeros comandos}
\begin{lstlisting}[style=rstyle]
2 + 3
sqrt(2)
sin(pi / 4)
log(exp(1))

x <- 5
y <- 2 * x + 1
y
\end{lstlisting}

\section{Ayuda integrada}
La ayuda de R es una herramienta de trabajo, no un recurso secundario. Algunas formas útiles de consultarla son:

\begin{lstlisting}[style=rstyle]
?sqrt
help(mean)
example(plot)
apropos("linear")
\end{lstlisting}

\begin{aviso}
Aprender a leer documentación técnica forma parte de la asignatura. En las prácticas se debe consultar la ayuda de R antes de buscar soluciones externas.
\end{aviso}

\chapter{Objetos, tipos de datos y expresiones}

\section{Objetos y asignación}
Un objeto es un nombre asociado a un valor. En R se recomienda usar \texttt{<-} para asignar.

\begin{lstlisting}[style=rstyle]
radio <- 3
area <- pi * radio^2
area
\end{lstlisting}

Los nombres deben ser descriptivos. \texttt{area\_circulo} es preferible a \texttt{a} cuando el contexto no es inmediato.

\section{Tipos básicos}
\begin{lstlisting}[style=rstyle]
entero <- 4L
real <- 4.5
texto <- "Laboratorio"
logico <- TRUE

class(entero)
class(real)
class(texto)
class(logico)
\end{lstlisting}

\section{Errores de redondeo}
Los números reales se representan de forma aproximada en el ordenador. Por ello, comparaciones aparentemente obvias pueden fallar.

\begin{lstlisting}[style=rstyle]
0.1 + 0.2 == 0.3
all.equal(0.1 + 0.2, 0.3)
\end{lstlisting}

En cálculo computacional conviene distinguir igualdad matemática exacta e igualdad numérica aproximada.

\chapter{Funciones}

\section{Funciones predefinidas}
Una función recibe argumentos y devuelve un resultado. Algunas funciones importantes en las primeras semanas son:

\begin{multicols}{2}
\begin{itemize}
  \item \texttt{sqrt}, \texttt{abs}, \texttt{round}
  \item \texttt{sin}, \texttt{cos}, \texttt{tan}
  \item \texttt{exp}, \texttt{log}
  \item \texttt{sum}, \texttt{prod}
  \item \texttt{mean}, \texttt{sd}, \texttt{var}
  \item \texttt{length}, \texttt{seq}, \texttt{rep}
\end{itemize}
\end{multicols}

\section{Funciones propias}
\begin{lstlisting}[style=rstyle]
area_circulo <- function(r) {
  pi * r^2
}

area_circulo(1)
area_circulo(3)
\end{lstlisting}

Una función debe resolver una tarea concreta. Si una función hace demasiadas cosas, suele ser más difícil de probar y reutilizar.

\section{Ejemplo: distancia euclídea}
\begin{lstlisting}[style=rstyle]
distancia <- function(x1, y1, x2, y2) {
  sqrt((x2 - x1)^2 + (y2 - y1)^2)
}

distancia(0, 0, 3, 4)
\end{lstlisting}

\chapter{Control de flujo}

\section{Condicionales}
Los condicionales permiten que un programa tome decisiones.

\begin{lstlisting}[style=rstyle]
clasificar_numero <- function(x) {
  if (x > 0) {
    "positivo"
  } else if (x < 0) {
    "negativo"
  } else {
    "cero"
  }
}

clasificar_numero(-2)
\end{lstlisting}

\section{Bucles}
Un bucle repite una operación. En R se usan bucles \texttt{for} y \texttt{while}; también existen alternativas vectorizadas.

\begin{lstlisting}[style=rstyle]
suma_cuadrados <- 0

for (k in 1:10) {
  suma_cuadrados <- suma_cuadrados + k^2
}

suma_cuadrados
\end{lstlisting}

\section{Ejemplo: algoritmo de Euclides}
\begin{lstlisting}[style=rstyle]
mcd <- function(a, b) {
  a <- abs(a)
  b <- abs(b)
  while (b != 0) {
    resto <- a %% b
    a <- b
    b <- resto
  }
  a
}

mcd(84, 30)
\end{lstlisting}

\chapter{Vectores, matrices y estructuras de datos}

\section{Vectores}
El vector es una estructura central en R. Muchas operaciones se aplican elemento a elemento.

\begin{lstlisting}[style=rstyle]
x <- c(1, 2, 3, 4, 5)
x^2
sum(x)
mean(x)
x[x > 3]
\end{lstlisting}

\section{Secuencias}
\begin{lstlisting}[style=rstyle]
seq(0, 1, by = 0.1)
seq(0, 2*pi, length.out = 100)
rep(1, times = 5)
\end{lstlisting}

\section{Matrices}
\begin{lstlisting}[style=rstyle]
A <- matrix(c(1, 2, 3, 4), nrow = 2, byrow = TRUE)
b <- c(1, 1)

A
t(A)
det(A)
solve(A, b)
\end{lstlisting}

\section{Tablas de datos}
Un \texttt{data.frame} organiza variables en columnas y observaciones en filas.

\begin{lstlisting}[style=rstyle]
datos <- data.frame(
  estudiante = c("A", "B", "C", "D"),
  practica = c(7.5, 8.0, 6.5, 9.0),
  examen = c(6.0, 7.0, 6.0, 8.5)
)

datos$media <- 0.5 * datos$practica + 0.5 * datos$examen
datos
\end{lstlisting}

\chapter{Representación gráfica}

\section{Gráficos base}
\begin{lstlisting}[style=rstyle]
x <- seq(-2*pi, 2*pi, length.out = 400)
y <- sin(x)

plot(x, y, type = "l",
     main = "Función seno",
     xlab = "x", ylab = "sin(x)")
abline(h = 0, col = "gray")
\end{lstlisting}

\section{\texttt{ggplot2}}
El paquete \texttt{ggplot2} permite construir gráficos por capas. La idea principal es separar datos, estética y geometría.

\begin{lstlisting}[style=rstyle]
# install.packages("ggplot2")
library(ggplot2)

datos <- data.frame(
  x = seq(-3, 3, length.out = 200)
)
datos$y <- dnorm(datos$x)

ggplot(datos, aes(x = x, y = y)) +
  geom_line(color = "steelblue", linewidth = 1) +
  labs(title = "Densidad normal estándar",
       x = "x", y = "f(x)") +
  theme_minimal()
\end{lstlisting}

\section{Criterios para un buen gráfico}
\begin{itemize}
  \item Debe responder a una pregunta concreta.
  \item Los ejes deben estar rotulados.
  \item Las unidades deben indicarse cuando existan.
  \item El título debe describir el contenido, no decorar.
  \item El gráfico debe poder interpretarse sin leer el código.
\end{itemize}

\chapter{Cálculo numérico elemental}

\section{Evaluación de funciones}
El ordenador permite evaluar funciones en mallas de puntos y aproximar propiedades globales.

\begin{lstlisting}[style=rstyle]
f <- function(x) x^3 - x - 1
x <- seq(0, 2, length.out = 200)
y <- f(x)
plot(x, y, type = "l")
abline(h = 0, col = "gray")
\end{lstlisting}

\section{Búsqueda de raíces}
R incluye \texttt{uniroot} para encontrar una raíz en un intervalo donde la función cambia de signo.

\begin{lstlisting}[style=rstyle]
f <- function(x) x^3 - x - 1
uniroot(f, interval = c(1, 2))
\end{lstlisting}

\section{Derivación numérica}
Una aproximación simple de la derivada es el cociente incremental centrado:

\[
f'(x) \approx \frac{f(x+h)-f(x-h)}{2h}.
\]

\begin{lstlisting}[style=rstyle]
derivada_centrada <- function(f, x, h = 1e-5) {
  (f(x + h) - f(x - h)) / (2 * h)
}

derivada_centrada(sin, pi / 4)
cos(pi / 4)
\end{lstlisting}

\section{Integración numérica}
La regla del trapecio aproxima el área bajo una curva.

\begin{lstlisting}[style=rstyle]
trapecio <- function(f, a, b, n = 1000) {
  x <- seq(a, b, length.out = n + 1)
  y <- f(x)
  h <- (b - a) / n
  h * (sum(y) - (y[1] + y[n + 1]) / 2)
}

trapecio(sin, 0, pi)
\end{lstlisting}

\chapter{Simulación}

\section{Números pseudoaleatorios}
La simulación usa números pseudoaleatorios. Para que un resultado sea reproducible se fija una semilla.

\begin{lstlisting}[style=rstyle]
set.seed(123)
runif(5)
rnorm(5)
\end{lstlisting}

\section{Monte Carlo para estimar \texorpdfstring{$\pi$}{pi}}
Una forma clásica de estimar \(\pi\) consiste en lanzar puntos al cuadrado \([-1,1]\times[-1,1]\) y contar cuántos caen dentro del círculo unidad.

\begin{lstlisting}[style=rstyle]
set.seed(1)
n <- 10000
x <- runif(n, -1, 1)
y <- runif(n, -1, 1)
dentro <- x^2 + y^2 <= 1
pi_estimado <- 4 * mean(dentro)
pi_estimado
\end{lstlisting}

\section{Error de simulación}
La precisión de una estimación de Monte Carlo mejora lentamente. Duplicar la precisión suele requerir multiplicar el número de simulaciones por cuatro.

\chapter{Tratamiento de datos y estadística inicial}

\section{Importación de datos}
El trabajo con datos requiere distinguir el archivo original, el código de limpieza y la tabla final. Los datos originales no deben modificarse manualmente.

\begin{lstlisting}[style=rstyle]
# datos <- read.csv("datos/observaciones.csv")
\end{lstlisting}

\section{Resumen descriptivo}
\begin{lstlisting}[style=rstyle]
set.seed(42)
notas <- data.frame(
  grupo = rep(c("A", "B"), each = 20),
  calificacion = c(rnorm(20, 6.5, 1), rnorm(20, 7.2, 1.2))
)

aggregate(calificacion ~ grupo, data = notas, mean)
aggregate(calificacion ~ grupo, data = notas, sd)
\end{lstlisting}

\section{Regresión lineal}
\begin{lstlisting}[style=rstyle]
set.seed(10)
x <- 1:30
y <- 2 + 0.7 * x + rnorm(30, sd = 2)

modelo <- lm(y ~ x)
summary(modelo)
plot(x, y)
abline(modelo, col = "red")
\end{lstlisting}

La salida de un modelo no debe copiarse sin interpretación. En un informe conviene explicar qué representa cada coeficiente, qué variabilidad queda sin explicar y qué limitaciones tiene el modelo.

\chapter{Edición científica con LaTeX}

\section{Estructura mínima}
\begin{lstlisting}[style=texstyle]
\documentclass{article}
\usepackage[spanish]{babel}
\usepackage{amsmath,amssymb}

\title{Mi primer documento}
\author{Nombre del estudiante}
\date{\today}

\begin{document}
\maketitle

\section{Introducción}
Este documento contiene texto y fórmulas.

\[
  \int_0^1 x^2\,dx = \frac{1}{3}.
\]

\end{document}
\end{lstlisting}

\section{Fórmulas}
LaTeX separa las fórmulas en línea, como \(a^2+b^2=c^2\), de las fórmulas destacadas:

\[
\sum_{k=1}^{n} k = \frac{n(n+1)}{2}.
\]

\section{Tablas}
\begin{lstlisting}[style=texstyle]
\begin{tabular}{lrr}
\toprule
Método & Error & Tiempo \\
\midrule
A & 0.013 & 0.42 \\
B & 0.008 & 0.71 \\
\bottomrule
\end{tabular}
\end{lstlisting}

\section{Buenas prácticas}
\begin{itemize}
  \item Separar contenido, estructura y formato.
  \item Usar etiquetas y referencias cruzadas para ecuaciones, figuras y tablas.
  \item Evitar capturas de pantalla de código o resultados cuando pueden incluirse como texto.
  \item Compilar con frecuencia para detectar errores pronto.
\end{itemize}

\chapter{Markdown, R Markdown y Quarto}

\section{Markdown}
Markdown permite escribir texto estructurado con una sintaxis ligera.

\begin{lstlisting}[style=plainstyle]
# Título

Un párrafo con **negrita** y una fórmula en línea: $x^2$.

- Primer elemento
- Segundo elemento
\end{lstlisting}

\section{Cuadernos reproducibles}
Un cuaderno reproducible combina texto, código y resultados. La idea central es que el informe se genere desde el código, no que el estudiante copie manualmente salidas de consola.

\begin{lstlisting}[style=plainstyle]
---
title: "Informe de práctica"
format: pdf
---

## Cálculo

```{r}
x <- seq(0, 1, length.out = 100)
mean(x^2)
```
\end{lstlisting}

\section{Estructura recomendada de un informe}
\begin{enumerate}
  \item Pregunta o problema.
  \item Método computacional.
  \item Código esencial.
  \item Resultados.
  \item Interpretación.
  \item Limitaciones o comprobaciones.
\end{enumerate}

\part{Prácticas de laboratorio}

\chapter{Guía general de prácticas}

\section{Formato común de entrega}
Cada práctica debe entregarse como carpeta comprimida o repositorio con:
\begin{itemize}
  \item un script principal \texttt{practicaXX.R};
  \item los datos usados, si procede;
  \item un informe breve en PDF, HTML o Quarto;
  \item las figuras generadas por el código, si no están incrustadas en el informe.
\end{itemize}

\section{Rúbrica común}
\begin{longtable}{p{0.30\textwidth}p{0.12\textwidth}p{0.48\textwidth}}
\toprule
Criterio & Peso & Evidencia esperada \\
\midrule
Corrección matemática & 30\% & Resultados coherentes, comprobaciones y uso adecuado de conceptos. \\
Código & 30\% & Claridad, nombres descriptivos, funciones cuando proceda, ausencia de duplicación innecesaria. \\
Reproducibilidad & 20\% & Proyecto organizado, rutas relativas, semilla aleatoria cuando proceda, informe generado desde código. \\
Comunicación & 20\% & Explicación breve, gráficos legibles, interpretación de resultados. \\
\bottomrule
\end{longtable}

\chapter{Práctica 1. Entorno de trabajo y primer proyecto}

\begin{objetivos}
\begin{itemize}
  \item Crear un proyecto de RStudio/Posit.
  \item Organizar carpetas de trabajo.
  \item Ejecutar comandos básicos en consola y script.
  \item Guardar un primer informe de resultados.
\end{itemize}
\end{objetivos}

\section{Actividad}
\begin{enumerate}
  \item Crear una carpeta de proyecto con subcarpetas \texttt{datos}, \texttt{scripts}, \texttt{figuras} e \texttt{informes}.
  \item Escribir un script que calcule áreas y volúmenes de figuras elementales.
  \item Incluir al menos una función propia.
  \item Guardar los resultados principales en un archivo de texto o informe breve.
\end{enumerate}

\begin{lstlisting}[style=rstyle]
area_rectangulo <- function(base, altura) {
  base * altura
}

volumen_esfera <- function(radio) {
  4 / 3 * pi * radio^3
}

area_rectangulo(3, 5)
volumen_esfera(2)
\end{lstlisting}

\begin{entrega}
Entregar el proyecto organizado y un informe de una página con los comandos usados, resultados y una reflexión sobre la estructura de carpetas.
\end{entrega}

\chapter{Práctica 2. Cálculo elemental y precisión numérica}

\begin{objetivos}
\begin{itemize}
  \item Trabajar con expresiones numéricas.
  \item Reconocer errores de redondeo.
  \item Comparar resultados exactos y aproximados.
\end{itemize}
\end{objetivos}

\section{Actividad}
\begin{enumerate}
  \item Evaluar expresiones con potencias, raíces, logaritmos y trigonometría.
  \item Comprobar identidades matemáticas con valores numéricos.
  \item Investigar casos donde la comparación exacta falla por precisión finita.
  \item Escribir una función \texttt{casi\_igual(x, y, tol)}.
\end{enumerate}

\begin{lstlisting}[style=rstyle]
casi_igual <- function(x, y, tol = 1e-8) {
  abs(x - y) < tol
}

casi_igual(0.1 + 0.2, 0.3)
\end{lstlisting}

\begin{entrega}
Entregar un script y un informe con tres ejemplos de identidades verificadas numéricamente y una explicación de las discrepancias observadas.
\end{entrega}

\chapter{Práctica 3. Funciones propias}

\begin{objetivos}
\begin{itemize}
  \item Diseñar funciones con argumentos claros.
  \item Probar funciones con casos simples.
  \item Documentar el comportamiento esperado.
\end{itemize}
\end{objetivos}

\section{Actividad}
Programar funciones para:
\begin{enumerate}
  \item convertir grados a radianes y radianes a grados;
  \item calcular la distancia entre dos puntos del plano;
  \item calcular el valor de un polinomio en un punto;
  \item normalizar un vector dividiendo por su norma euclídea.
\end{enumerate}

\begin{lstlisting}[style=rstyle]
evaluar_polinomio <- function(coeficientes, x) {
  potencias <- 0:(length(coeficientes) - 1)
  sum(coeficientes * x^potencias)
}

evaluar_polinomio(c(1, 0, 3), 2)
\end{lstlisting}

\begin{entrega}
Entregar código con pruebas de cada función y un breve comentario sobre posibles entradas no válidas.
\end{entrega}

\chapter{Práctica 4. Condicionales, bucles y algoritmos elementales}

\begin{objetivos}
\begin{itemize}
  \item Implementar algoritmos con control de flujo.
  \item Distinguir bucles \texttt{for} y \texttt{while}.
  \item Incorporar comprobaciones de entrada.
\end{itemize}
\end{objetivos}

\section{Actividad}
\begin{enumerate}
  \item Implementar el algoritmo de Euclides para el máximo común divisor.
  \item Programar una función que determine si un entero es primo.
  \item Calcular los primeros \(n\) términos de la sucesión de Fibonacci.
  \item Comparar el tiempo de ejecución para distintos tamaños de entrada.
\end{enumerate}

\begin{lstlisting}[style=rstyle]
es_primo <- function(n) {
  if (n < 2) return(FALSE)
  if (n == 2) return(TRUE)
  for (k in 2:floor(sqrt(n))) {
    if (n %% k == 0) return(FALSE)
  }
  TRUE
}
\end{lstlisting}

\begin{entrega}
Entregar funciones, pruebas y una tabla con resultados para varios valores de \(n\).
\end{entrega}

\chapter{Práctica 5. Vectores y matrices}

\begin{objetivos}
\begin{itemize}
  \item Usar operaciones vectorizadas.
  \item Resolver sistemas lineales pequeños.
  \item Interpretar resultados matriciales.
\end{itemize}
\end{objetivos}

\section{Actividad}
\begin{enumerate}
  \item Crear vectores mediante \texttt{c}, \texttt{seq} y \texttt{rep}.
  \item Calcular normas, productos escalares y distancias.
  \item Construir matrices y operar con transpuestas, determinantes e inversas.
  \item Resolver sistemas \(Ax=b\) con \texttt{solve}.
\end{enumerate}

\begin{lstlisting}[style=rstyle]
A <- matrix(c(2, 1, 1, 3), nrow = 2)
b <- c(1, 4)
x <- solve(A, b)
A %*% x
\end{lstlisting}

\begin{entrega}
Entregar un informe con al menos tres sistemas lineales y una verificación \(Ax=b\) para cada uno.
\end{entrega}

\chapter{Práctica 6. Datos en tablas}

\begin{objetivos}
\begin{itemize}
  \item Crear y manipular \texttt{data.frame}.
  \item Calcular nuevas columnas.
  \item Filtrar observaciones.
\end{itemize}
\end{objetivos}

\section{Actividad}
\begin{enumerate}
  \item Crear una tabla con datos simulados de estudiantes y calificaciones.
  \item Calcular nota final ponderada.
  \item Clasificar a los estudiantes como apto/no apto.
  \item Obtener medias por grupo.
\end{enumerate}

\begin{lstlisting}[style=rstyle]
set.seed(5)
datos <- data.frame(
  grupo = rep(c("A", "B", "C"), each = 10),
  practica = runif(30, 4, 10),
  examen = runif(30, 3, 10)
)

datos$final <- 0.5 * datos$practica + 0.5 * datos$examen
datos$apto <- datos$final >= 5
aggregate(final ~ grupo, data = datos, mean)
\end{lstlisting}

\begin{entrega}
Entregar tabla final, código y comentario sobre la distribución de calificaciones por grupo.
\end{entrega}

\chapter{Práctica 7. Visualización de funciones y datos}

\begin{objetivos}
\begin{itemize}
  \item Representar funciones matemáticas.
  \item Construir gráficos de dispersión e histogramas.
  \item Mejorar títulos, ejes y leyendas.
\end{itemize}
\end{objetivos}

\section{Actividad}
\begin{enumerate}
  \item Representar \(f(x)=\sin(x)\), \(g(x)=\cos(x)\) y \(h(x)=e^{-x^2}\).
  \item Crear un histograma de una muestra normal.
  \item Crear un diagrama de dispersión con recta de ajuste.
  \item Exportar al menos una figura a la carpeta \texttt{figuras}.
\end{enumerate}

\begin{lstlisting}[style=rstyle]
png("figuras/seno_coseno.png", width = 900, height = 600)
x <- seq(-2*pi, 2*pi, length.out = 400)
plot(x, sin(x), type = "l", col = "blue", ylab = "valor")
lines(x, cos(x), col = "red")
legend("topright", legend = c("sen(x)", "cos(x)"),
       col = c("blue", "red"), lty = 1)
dev.off()
\end{lstlisting}

\begin{entrega}
Entregar script, figuras generadas y comentario sobre qué gráfico comunica mejor cada tipo de información.
\end{entrega}

\chapter{Práctica 8. Raíces y aproximación numérica}

\begin{objetivos}
\begin{itemize}
  \item Resolver ecuaciones no lineales sencillas.
  \item Usar \texttt{uniroot}.
  \item Comparar aproximaciones con soluciones conocidas.
\end{itemize}
\end{objetivos}

\section{Actividad}
\begin{enumerate}
  \item Representar \(f(x)=x^3-x-1\) en varios intervalos.
  \item Localizar un intervalo con cambio de signo.
  \item Calcular una raíz con \texttt{uniroot}.
  \item Implementar el método de bisección como función propia.
\end{enumerate}

\begin{lstlisting}[style=rstyle]
biseccion <- function(f, a, b, tol = 1e-8, max_iter = 100) {
  if (f(a) * f(b) > 0) stop("No hay cambio de signo")
  for (i in 1:max_iter) {
    c <- (a + b) / 2
    if (abs(f(c)) < tol || (b - a) / 2 < tol) return(c)
    if (f(a) * f(c) < 0) {
      b <- c
    } else {
      a <- c
    }
  }
  (a + b) / 2
}
\end{lstlisting}

\begin{entrega}
Comparar \texttt{biseccion} y \texttt{uniroot} en tres funciones distintas.
\end{entrega}

\chapter{Práctica 9. Derivación e integración numérica}

\begin{objetivos}
\begin{itemize}
  \item Aproximar derivadas mediante diferencias finitas.
  \item Aproximar integrales mediante reglas de cuadratura.
  \item Analizar el efecto del tamaño de paso.
\end{itemize}
\end{objetivos}

\section{Actividad}
\begin{enumerate}
  \item Implementar diferencias hacia delante y centradas.
  \item Aproximar derivadas de \(\sin(x)\), \(e^x\) y \(x^3\).
  \item Implementar la regla del trapecio.
  \item Estudiar cómo cambia el error al variar \(h\) o \(n\).
\end{enumerate}

\begin{entrega}
Entregar una tabla de errores y un gráfico que muestre la relación entre precisión y tamaño de paso.
\end{entrega}

\chapter{Práctica 10. Simulación y Monte Carlo}

\begin{objetivos}
\begin{itemize}
  \item Fijar semillas aleatorias.
  \item Simular experimentos.
  \item Estimar probabilidades y magnitudes geométricas.
\end{itemize}
\end{objetivos}

\section{Actividad}
\begin{enumerate}
  \item Simular lanzamientos de una moneda equilibrada.
  \item Simular tiradas de dos dados y estimar la probabilidad de suma 7.
  \item Estimar \(\pi\) mediante Monte Carlo.
  \item Repetir el experimento con distintos tamaños muestrales.
\end{enumerate}

\begin{entrega}
Entregar resultados para al menos cuatro tamaños muestrales y explicar la variabilidad observada.
\end{entrega}

\chapter{Práctica 11. Estadística descriptiva}

\begin{objetivos}
\begin{itemize}
  \item Calcular medidas de posición y dispersión.
  \item Visualizar distribuciones.
  \item Comparar grupos.
\end{itemize}
\end{objetivos}

\section{Actividad}
\begin{enumerate}
  \item Generar o importar una tabla de datos con una variable cuantitativa y una variable de grupo.
  \item Calcular media, mediana, desviación típica y rango intercuartílico.
  \item Construir histogramas y diagramas de caja.
  \item Redactar una comparación de grupos.
\end{enumerate}

\begin{entrega}
Entregar un informe con tablas, gráficos y una interpretación de máximo dos páginas.
\end{entrega}

\chapter{Práctica 12. Regresión lineal}

\begin{objetivos}
\begin{itemize}
  \item Ajustar un modelo lineal simple.
  \item Interpretar pendiente e intercepto.
  \item Diagnosticar visualmente el ajuste.
\end{itemize}
\end{objetivos}

\section{Actividad}
\begin{enumerate}
  \item Simular datos con relación lineal y ruido.
  \item Ajustar \texttt{lm(y \textasciitilde{} x)}.
  \item Dibujar la nube de puntos y la recta ajustada.
  \item Analizar residuos mediante gráficos.
\end{enumerate}

\begin{lstlisting}[style=rstyle]
set.seed(20)
x <- runif(50, 0, 10)
y <- 1.5 + 2.2 * x + rnorm(50, sd = 2)
modelo <- lm(y ~ x)
summary(modelo)
plot(modelo)
\end{lstlisting}

\begin{entrega}
Entregar informe reproducible con interpretación de los coeficientes y comentario sobre la calidad del ajuste.
\end{entrega}

\chapter{Práctica 13. Documento científico en LaTeX}

\begin{objetivos}
\begin{itemize}
  \item Crear un documento LaTeX desde cero.
  \item Escribir fórmulas, tablas y referencias cruzadas.
  \item Compilar un PDF correctamente.
\end{itemize}
\end{objetivos}

\section{Actividad}
\begin{enumerate}
  \item Crear un documento con título, autor, resumen y dos secciones.
  \item Incluir al menos tres fórmulas matemáticas.
  \item Incluir una tabla con resultados numéricos.
  \item Usar etiquetas y referencias cruzadas.
  \item Añadir una bibliografía mínima.
\end{enumerate}

\begin{entrega}
Entregar el archivo \texttt{.tex}, el PDF compilado y cualquier archivo auxiliar necesario para la bibliografía.
\end{entrega}

\chapter{Práctica 14. Informe reproducible con R Markdown o Quarto}

\begin{objetivos}
\begin{itemize}
  \item Combinar texto, código y resultados.
  \item Generar un informe reproducible.
  \item Evitar copiar manualmente resultados.
\end{itemize}
\end{objetivos}

\section{Actividad}
\begin{enumerate}
  \item Crear un archivo \texttt{.qmd} o \texttt{.Rmd}.
  \item Incluir una introducción al problema.
  \item Ejecutar código R dentro del documento.
  \item Mostrar al menos un gráfico generado desde el propio informe.
  \item Renderizar a HTML o PDF.
\end{enumerate}

\begin{entrega}
Entregar el archivo fuente y el informe renderizado. El informe debe poder regenerarse sin cambios manuales.
\end{entrega}

\chapter{Práctica 15. Proyecto integrador}

\begin{objetivos}
\begin{itemize}
  \item Integrar programación, datos, gráficos, cálculo y comunicación científica.
  \item Trabajar con una estructura de proyecto completa.
  \item Presentar resultados de forma defendible.
\end{itemize}
\end{objetivos}

\section{Propuesta}
El estudiante elegirá uno de los siguientes proyectos:
\begin{enumerate}
  \item Estudio numérico de una familia de funciones y sus raíces.
  \item Simulación de un experimento aleatorio y estimación de probabilidades.
  \item Análisis descriptivo y gráfico de un conjunto de datos.
  \item Ajuste e interpretación de un modelo lineal simple.
\end{enumerate}

\section{Requisitos mínimos}
\begin{itemize}
  \item Código R organizado en al menos un script principal.
  \item Al menos una función propia.
  \item Al menos dos gráficos interpretados.
  \item Informe final en LaTeX, R Markdown o Quarto.
  \item Conclusiones que conecten los resultados con el problema planteado.
\end{itemize}

\begin{entrega}
Entregar proyecto completo y realizar una defensa breve de 5 minutos. La defensa debe explicar el problema, el método, los resultados y una limitación del trabajo.
\end{entrega}

\appendix

\chapter{Referencia rápida de R}

\section{Objetos y operaciones}
\begin{lstlisting}[style=rstyle]
x <- 3
y <- c(1, 2, 3)
z <- seq(0, 1, length.out = 5)
sum(y)
mean(y)
length(y)
\end{lstlisting}

\section{Funciones}
\begin{lstlisting}[style=rstyle]
nombre_funcion <- function(argumento1, argumento2) {
  resultado <- argumento1 + argumento2
  resultado
}
\end{lstlisting}

\section{Control de flujo}
\begin{lstlisting}[style=rstyle]
if (condicion) {
  # instrucciones
} else {
  # instrucciones alternativas
}

for (k in 1:10) {
  print(k)
}
\end{lstlisting}

\chapter{Plantilla de informe de práctica}

\begin{lstlisting}[style=plainstyle]
Título de la práctica
Nombre:
Fecha:

1. Problema
Descripción breve del problema que se quiere resolver.

2. Método
Explicación de las funciones, algoritmos o datos usados.

3. Código esencial
Fragmentos principales o referencia al script.

4. Resultados
Tablas, valores numéricos y gráficos.

5. Interpretación
Qué significan los resultados y qué limitaciones tienen.
\end{lstlisting}

\chapter{Bibliografía inicial}
\begin{itemize}
  \item Matloff, N. \emph{The Art of R Programming}. No Starch Press.
  \item Wickham, H. \emph{Advanced R}. Chapman and Hall/CRC.
  \item Wickham, H. y Grolemund, G. \emph{R for Data Science}. O'Reilly Media.
  \item Grätzer, G. \emph{More Math Into LaTeX}. Springer.
  \item Xie, Y., Dervieux, C. y Riederer, E. \emph{R Markdown Cookbook}. Chapman and Hall/CRC.
  \item Venables, W. N., Smith, D. M. y R Core Team. \emph{An Introduction to R}.
  \item Lamport, L. \emph{LaTeX: A Document Preparation System}. Addison-Wesley.
  \item Xie, Y., Allaire, J. J. y Grolemund, G. \emph{R Markdown: The Definitive Guide}. Chapman and Hall/CRC.
\end{itemize}

\backmatter

\chapter*{Notas para revisión docente}
Este borrador se ha preparado a partir de la guía docente disponible. Antes de usarlo como manual definitivo conviene revisar:
\begin{itemize}
  \item calendario académico real y festivos;
  \item criterios exactos de evaluación aprobados por la coordinación;
  \item software disponible en aulas;
  \item conjuntos de datos oficiales de la asignatura;
  \item formato institucional exigido para guiones de prácticas;
  \item nivel inicial real del grupo.
\end{itemize}

\end{document}
